\subsubsection{7.1 Prompt Injection and LLM-Specific
Defences}\label{prompt-injection-and-llm-specific-defences}

The original prompt-injection taxonomy is due to Perez and Ribeiro
{[}@perez2022ignore{]}; the realistic indirect-injection threat model
originated with Greshake et al. {[}@greshake2023not{]}. Subsequent
defences fall into three families. \emph{Input perturbation} defences
--- Spotlighting {[}@hines2024defending{]}, paraphrasing, retokenisation
--- modify the agent's context so that untrusted regions are visually
marked or syntactically distinguished. \emph{Output filtering} defences
--- constitutional classifiers {[}@anthropic2024constitutional{]},
commercial prompt shields --- apply a secondary model to the agent's
output to catch policy violations. \emph{Training-time defences} ---
StruQ {[}@chen2024struq{]}, SecAlign {[}@chen2025secalign{]},
instruction-hierarchy fine-tuning --- modify the base model. All three
families share the architectural property that they operate over
unbounded natural language and provide statistical rather than
structural guarantees. The empirical ceiling on these approaches against
modern benchmarks sits around 14-31\% attack success, as Section 6
reports.

The benchmark literature --- AgentDojo {[}@debenedetti2024agentdojo{]},
InjecAgent {[}@zhan2024injecagent{]}, TensorTrust, AdvBench --- has been
instrumental in disciplining defence claims. We adopt AgentDojo and
InjecAgent as our headline benchmarks because they specifically measure
tool-call-mediated success, which is the class VCD addresses.

\textbf{A note on policy-targeted text scanning versus generic PI
detection.} In §6.2 the text-only configuration of VCD
(\texttt{vcd\_text}) outperforms prompt-shields' DeBERTa detector on the
less capable v3.2 model --- 0.7\% ASR with 79.9\% benign retention,
versus shields' 7.6\% ASR at 33.3\% benign. We did not design
\texttt{vcd\_text} to compete head-on with shields and we did not
anticipate it would dominate; we offer the following mechanism. A
generic PI detector classifies natural-language prose against patterns
of \emph{injection-style} writing. The classifier is policy-agnostic ---
it does not know what the agent is supposed to be allowed to do, and
consequently must err toward conservatism on any text whose stylistic
markers resemble its training distribution. On a less-capable base model
whose reasoning trajectory is itself noisier, the detector's
false-positive rate compounds: many legitimate user requests, expressed
in the way real users write, look stylistically similar to injection
prose. VCD's text scanner, by contrast, is policy-aware --- it carries
the in-force capability's constraint set, and flags only content that
would steer the agent toward calls those constraints disallow. The
surface-form of the injection prose is irrelevant; only the structural
pattern of ``this would induce a gate-disallowed call'' matters. On a
more capable model, both detectors converge because the LLM's own
discrimination eats up most of the marginal precision; on a less capable
model, the policy-aware approach pays its full dividend. The argument is
in line with the general observation that, where structure exists,
structural enforcement strictly dominates statistical detection.

Greshake et al.'s indirect-injection threat model {[}greshake2023not{]}
is the closest framing to ours --- prompt injection as a
privilege-escalation problem reachable via untrusted documents in the
agent's context --- and the recent thread of work on ``spotlighting +
tool authentication'' combinations explores adjacent ideas informally.
None of this prior work provides cryptographic enforcement, attenuation
invariants, or SMT-verified policy completeness, and to our knowledge
none has been evaluated against AgentDojo or InjecAgent with a
structurally-enforced gate.

\subsubsection{7.2 Capability Discipline}\label{capability-discipline}

Object capabilities were introduced by Levy {[}@levy1984capability{]}
and developed by Mark Miller in his dissertation on robust composition
{[}@miller2006robust{]}; major implementations include EROS
{[}@shapiro1999eros{]}, KeyKOS, Caja, and seL4 {[}@klein2009sel4{]}.
Macaroons {[}@birgisson2014macaroons{]} provide a particularly close
analogue to our attenuation primitive, with HMAC-based chained
restriction. The technical deltas are summarised in Table 1.

\textbf{Table 1. Macaroons vs VCD capability tokens.}

\begin{longtable}[]{@{}
  >{\raggedright\arraybackslash}p{(\linewidth - 4\tabcolsep) * \real{0.3333}}
  >{\raggedright\arraybackslash}p{(\linewidth - 4\tabcolsep) * \real{0.3333}}
  >{\raggedright\arraybackslash}p{(\linewidth - 4\tabcolsep) * \real{0.3333}}@{}}
\toprule\noalign{}
\begin{minipage}[b]{\linewidth}\raggedright
Property
\end{minipage} & \begin{minipage}[b]{\linewidth}\raggedright
Macaroons
\end{minipage} & \begin{minipage}[b]{\linewidth}\raggedright
VCD
\end{minipage} \\
\midrule\noalign{}
\endhead
\bottomrule\noalign{}
\endlastfoot
Cryptographic substrate & HMAC-SHA256 chain & Ed25519 signature per
token \\
Verifiable without issuer secret & No (verifier holds HMAC key) & Yes
(verifier holds only the public key) \\
Identifier & bearer-token byte string & content-addressed SHA-256 of
canonical body \\
Reassignment of identifier & possible (no binding) & precluded (id
determines body) \\
Constraint vocabulary & opaque caveat strings & typed constraint kinds
with lattice meet \\
Mechanised attenuation soundness & no & yes (Lean 4, §4) \\
Binding to a verified policy artifact & no &
\texttt{policy\_proof\_hash} field \\
Audit-graph composition & informal & content-addressed; chains into
Merkle log \\
\end{longtable}

The deltas are not academic. The public-key-verifiable property means a
downstream tool implementation can independently verify a token without
sharing a secret with the issuer --- relevant in multi-vendor agent
deployments. The content-addressed identifier blocks a class of
token-substitution attacks that affects bearer-token systems generally.
The typed constraint vocabulary is what enables the SMT bridge in
Theorem 3.

The industry-familiar comparison is to OAuth 2.0 scopes. OAuth's
\texttt{scope=read:email\ scope=send:email} model expresses
\emph{permission-level} restrictions; VCD's constraints are
\emph{value-level} (the per-call \texttt{recipient}, \texttt{amount},
and so on are checked, not just the verb). Macaroons' caveats are
value-level but informal strings; VCD's are a typed vocabulary that
admits SMT verification. OAuth scopes are documented; VCD's attenuation
invariants are mechanically enforced in Lean. OAuth tokens carry no
proof of policy completeness; VCD tokens cite a SMT-verified
\texttt{policy\_proof\_hash}. The three frameworks (OAuth scopes,
Macaroons, VCD) lie on a progression of structural commitment: OAuth
scopes are a vocabulary, Macaroons add chained attenuation under HMAC,
VCD adds value-level constraints + mechanised soundness + verified
policy citation.

Capability-based approaches to LLM agents have been discussed informally
in community forums and as part of broader ``tool allowlist'' practices
in major agent frameworks, but to our knowledge no prior published work
provides the mechanised attenuation invariants, cryptographic gate
enforcement, and empirical evaluation on prompt-injection benchmarks
that VCD combines.

\textbf{CaMeL.} The closest contemporaneous published work is CaMeL
{[}@debenedetti2025camel{]}, which independently identifies
capability-style metadata as a route to structurally-enforced agent
safety. CaMeL splits the agent into a \emph{privileged} LLM that emits
restricted Python and a \emph{quarantined} LLM that processes untrusted
data, with capability metadata flowing through interpreter values; on
AgentDojo, CaMeL reports 77\% benign task completion at provable
security, against an unprotected baseline of 84\%. CaMeL and VCD share
the framing --- move enforcement out of the model into a
structurally-checked layer --- and on this framing both contribute. The
technical deltas are: (a) CaMeL's soundness is the dynamic semantics of
a custom Python interpreter; VCD's three soundness theorems are
mechanised in Lean 4. (b) CaMeL requires architectural surgery --- two
LLMs, a custom planner/quarantine split --- which is incompatible with
deployed single-LLM agents and with the Model Context Protocol; VCD is a
sidecar gate that drops into existing agents without changing the LLM
call. (c) CaMeL's capabilities are interpreter-local metadata that exist
only within a single agent's execution; VCD's capability tokens are
public-key-verifiable signed artefacts that travel across organisations,
clouds, and tool implementations. (d) CaMeL relies on the user to author
policies inline; VCD binds each policy to an SMT-verified
\texttt{policy\_proof\_hash} whose verification is independent of the
issuer. The two approaches are not mutually exclusive --- a CaMeL-style
dual-LLM architecture composed with VCD-style cryptographic
authorisation is strictly stronger than either alone --- but only VCD
produces a portable audit artefact, and only VCD's soundness is
machine-checked.

\textbf{Microsoft Agent Governance Toolkit.} A second contemporaneous
effort is the Microsoft Agent Governance Toolkit (AGT), released as
seven MIT-licensed packages on 2 April 2026 {[}@microsoft2026agt{]} ---
Agent OS (an in-process policy gate), AgentMesh (SPIFFE/Ed25519 agent
identity), Agent Runtime (rings 0--3 sandbox), Agent Hypervisor
(Merkle-chained audit log), an MCP security gateway, and SRE /
compliance tooling. AGT and VCD overlap on the abstract claim ---
structural enforcement at the tool-call boundary, cryptographic
identity, append-only audit --- and on the deployment surface, since
both target the Microsoft Agent Framework's \texttt{FunctionMiddleware}
hook. The differences are load-bearing for what each contributes to the
literature. (a) \textbf{Soundness substrate.} AGT validates policy
decisions via conformance-test suites; VCD ships three theorems
mechanised in Lean 4 with zero \texttt{sorry}s (§4), each externally
re-checkable by rebuilding the published Lean development. (b)
\textbf{Audit artefact.} AGT's audit chain is operator-rooted: the chain
root is signed by the deploying organisation's audit key, and verifiers
need an attestation that the operator's audit-key rotation was correctly
performed. VCD receipts are content-addressed signed artefacts naming,
by hash, the schema, the policy proof, and the cited Lean theorem id; a
verifier holding only the issuer's published material can confirm the
receipt without contacting the operator (§7.5). (c) \textbf{Policy
substrate.} AGT supports YAML, OPA/Rego, and Cedar policy languages
interpreted at runtime by the gate; VCD compiles a policy through an SMT
prover (Z3) and produces a \texttt{policy\_proof\_hash} that certifies
the policy holds over every string the schema admits, before any token
referencing it is minted (§3.1, §6.2.1). AGT does not currently document
a third-party Policy Decision Point plug-in contract; this paper's
reference implementation ships an Agent Framework middleware
(\texttt{docs/proposals/agent-framework-middleware.md}) that composes
alongside AGT in production deployments rather than displacing it. We
take the strategic position that AGT and VCD are complements: AGT
supplies a rich in-process default policy engine; VCD supplies the
formally-verified policy decisions and offline-verifiable audit
artefacts that high-assurance verticals require.

\subsubsection{7.3 SMT for Security
Policy}\label{smt-for-security-policy}

The use of SMT solvers for security-policy reasoning has substantial
prior art. Margrave {[}@fisler2005margrave{]} used decision procedures
for XACML access-control reasoning; subsequent work scaled the approach
to AWS IAM policies via Backes et al.'s Zelkova system
{[}backes2018semantic{]}, which encodes IAM access policies as SMT and
proves containment and reachability properties. We borrow the conceptual
move from this line of work --- encode the policy as a violation
predicate and ask the solver to refute it --- and apply it to the
specific case of JSON Schema-bounded tool-call grammars in LLM agent
settings.

\subsubsection{7.4 Software Supply-Chain
Attestation}\label{software-supply-chain-attestation}

The framing of ``every artefact content-addressed, every link signed''
derives from the software supply-chain attestation literature. Sigstore
{[}@newman2022sigstore{]} and SLSA articulate this design pattern for
build provenance; in-toto {[}@torres2019intoto{]} extends to multi-step
pipelines; TUF {[}@samuel2010survivable{]} establishes the underlying
compromise-resilience properties. VCD applies the same architectural
pattern --- content addresses, signed transitions, pinned roots --- to a
different domain (agent tool-call policy) and a different shape of
artefact (capability tokens rather than build attestations).

\subsubsection{7.5 Capability Receipts as an Audit
Primitive}\label{capability-receipts-as-an-audit-primitive}

The remaining adjacencies --- Microsoft's Prompt Shields, AWS Bedrock
policy controls for agents, Google's Model Armor, Lakera Guard / Cisco
AI Defense, Pillar Security, Anthropic's constitutional classifiers ---
operate at the \emph{defensive} layer: they reduce attack-success rates
statistically. None produces a portable artefact tied to a verified
policy. None offers an audit log a third party can independently verify.
The question we pose, and answer affirmatively, is: \emph{what would it
take to give a regulator, an internal audit team, or a downstream tool
implementation a cryptographic record of agent activity that they could
verify without trusting the agent's operator?}

VCD's capability receipt is that record. A receipt has the form

\texttt{receipt\ =\ (issuer\_cert,\ schema\_hash,\ policy\_proof\_hash,\ lean\_theorem\_id,\ attenuation\_chain,\ call\_args\_hash,\ decision,\ timestamp)}

signed under the issuer's Ed25519 key (Figure 4). Every field is
content-addressed: \texttt{schema\_hash} is the SHA-256 of the canonical
JSON Schema the prover ran against; \texttt{policy\_proof\_hash} is the
hash of the Z3 proof artefact (or counterexample);
\texttt{lean\_theorem\_id} names the mechanised soundness theorem in
\texttt{paper/lean/}; \texttt{attenuation\_chain} lists every parent
token's id, ensuring the receipt's authority chains to a pinned root;
\texttt{call\_args\_hash} is the SHA-256 of the canonical-serialised
tool-call arguments the gate inspected. A verifier holding only the
issuer's public key and an offline copy of the published Lean
development can independently confirm five properties: the signature is
valid under the pinned key; the schema hash matches the published
registry entry; the policy proof hash matches the published proof
artefact; the cited Lean theorem closes (rebuild the development
locally); the call argument hash satisfies the cited policy's
constraints. None of those checks requires contacting the issuer.

The artefact is the durable contribution. Spotlighting and StruQ make
the agent harder to fool; AWS Bedrock policy controls make the agent's
actions cheaper to inspect; CaMeL makes the agent's reasoning
structurally compartmented. \textbf{VCD is the only published system
whose output is a verifiable record of agent authority.} Three concrete
uses motivate the framing:

\begin{enumerate}
\def\labelenumi{\arabic{enumi}.}
\item
  \textbf{Cross-organisation audit (Figure 5).} A bank's FCA examiner
  reviews three months of receipts emitted by the bank's AI
  customer-service agent. The bank has published, signed and
  out-of-band, four artefacts: its issuer public key, its policy
  registry (each policy hashed and signed), its schema registry, and the
  URL of its Lean development. The examiner requests the receipt log
  from the bank (which need not be confidential --- every receipt is
  independently verifiable). For each receipt the examiner runs five
  offline checks: signature against the published key;
  \texttt{schema\_hash} against the registry;
  \texttt{policy\_proof\_hash} against the registry; rebuild the Lean
  development and confirm the cited theorem closes; check that
  \texttt{call\_args\_hash} satisfies the policy's constraint set. The
  audit's conclusion is a deterministic function of the published
  artefacts. \textbf{No on-premise inspection of the bank's AI
  infrastructure is required; no shared secret with the bank's vendor is
  involved; no information asymmetry between bank and regulator
  remains.}
\item
  \textbf{Tool-side independent enforcement.} A payments-network
  operator receives a \texttt{transfer\_funds} call from a third-party
  agent platform. Before executing, the operator's gate verifies the
  capability receipt: was this call authorised by a token signed by a
  pinned issuer the operator recognises? Does the cited schema match the
  operator's published function declaration? Does the cited policy proof
  actually constrain the call args the operator received? The operator
  depends on no other piece of the agent platform.
\item
  \textbf{Inter-platform delegation.} A travel-booking agent at one
  organisation delegates to a payment-processing agent at another. The
  travel agent attenuates its capability token to a tighter child for
  the payment agent. The payment agent's gate, holding only the travel
  organisation's published issuer key, verifies the child token's
  attenuation chain back to a recognised root. Authority crosses an
  organisational boundary without an out-of-band trust agreement.
\end{enumerate}

These three uses are presented as worked examples rather than measured
deployments. The relevant claim is that the artefact is a structural fit
for an audit pattern the existing prompt-injection-defence literature
does not address, and that the absence of such an artefact in adjacent
commercial products (Prompt Shields, Bedrock policy controls, Lakera
Guard) is not an oversight but a property of the design space: a
statistical classifier cannot produce a verifiable receipt because there
is no verifiable claim to cite. A structurally-enforced gate composed
with a content-addressed signing chain can.
